% ============================================================
% ANNEXES
% ============================================================
\chapter*{Annexes}
\pagestyle{styleGlobal}
\label{sec:annexes}
\markboth{Annexes}{Annexes}
\addcontentsline{toc}{chapter}{Annexes}

\vspace*{0.5cm}

\section*{Annexe A : Schéma de la base de données (Prisma)}

Le schéma ci-dessous représente la définition complète de la base de données de « SALISA », telle qu'elle a été implémentée avec l'outil Prisma ORM. Il traduit directement le diagramme de classes présenté dans la partie Analyse et Conception et sert de référence technique pour la génération des modèles de données dans l'application.

\begin{figure}[H]
	\centering
	\includegraphics[width=0.95\textwidth]{schema_prisma.png}
	\caption{Schéma Prisma de la base de données de « SALISA »}
	\label{fig:schema_prisma}
\end{figure}

\section*{Annexe B : Organigramme original d'Akieni}

La figure suivante présente l'organigramme original fourni par la structure d'accueil Akieni, tel qu'il nous a été transmis. Il a servi de base à la rédaction de la section 1.1.3 du présent document.

\begin{figure}[H]
	\centering
	\includegraphics[width=0.85\textwidth]{organigramme_original_akieni.jpg}
	\caption{Organigramme original d'Akieni (document interne)}
	\label{fig:organigramme_original}
\end{figure}